\section*{4 Upper bounds (author source pp.22:8--22:10)}
We show in this section that if $F$ is satisfiable, then $\RRef_{d,N}(F,s)$ has small depth-$d$ refutations:
\begin{originalresult}{Proposition 4.1}
If $F$ is a satisfiable CNF formula with $n$ variables and $m$ clauses, then there is a depth-$d$ LK refutation of $\RRef_{d,N}(F,s)$ of size
\[S=O\left(\left(N^{d+3}s^2n(m+s^2n^3)\right)^2\right).\]
In particular, if $m=O(s^2n^3)$, then $S=O(N^{2(d+3)}(sn)^8)$.
\end{originalresult}
\begin{proof}
We start with a small depth-2 LK tree-like refutation of $\RRef(F,s)$. This refutation will be such that after the substitution with Sipser functions, we get a depth-$(d+1)$ tree-like refutation of $\RRef_{d,N}(F,s)$, which in turn we can convert to a depth-$d$ DAG-like refutation of $\RRef_{d,N}(F,s)$ by Proposition2.2.
We write, for better readability, $A_1,\ldots,A_k\to B_1,\ldots,B_\ell$ instead of $\bar A_1\vee\cdots\vee\bar A_k\vee B_1\vee\cdots\vee B_\ell$.
Let $\alpha$ be an assignment that satisfies every clause of $F$. We set
\[T(u):=P[u]\to\bigvee_{i=1}^n D[u,i,\alpha(x_i)].\]
What $T(u)$ says is that if $P[u]$ is true, then $\alpha$ satisfies the $u$-th clause in the refutation $\Ref(F,s)$ describes.
Our refutation of $\RRef(F,s)$ consists of $s-1$ stages, starting with stage0. In the $u$-th stage, $T(1),\ldots,T(s-u)\to0$ will have been derived. Then we can use this formula, along with a derivation of $T(1),\ldots,T(s-u-1)\to T(s-u)$, to derive $T(1),\ldots,T(s-u-1)\to0$. In the $s-1$-th stage, $T(1)\to0$ will have been derived, at which point we can reach a contradiction by deriving $T(1)$.
A derivation of $T(1),\ldots,T(v-1)\to T(v)$ is sketched in Figure1.
\begin{figure}[ht]\centering\includegraphics[page=9,width=.82\textwidth,trim=90 400 100 210,clip]{frege-original.pdf}\caption{Original Figure1: A sketch of a derivation of $T(1),\ldots,T(v-1)\to T(v)$.}\end{figure}
The formulas $I[v,j]\to T(v)$ for $j\in[m]$ can be immediately derived from the clauses $P[u]\wedge I[v,j]\to D[v,i,\alpha(x_i)]$, which are clauses corresponding to condition(3.18), as the fact that $\alpha$ satisfies the clause $C_j$ means that $x_i^{\alpha(x_i)}$ must belong to $C_j$ for some $i$. These formulas can be in turn used along with the clauses(3.11) for $I[v,k]$ and(3.12) to derive $V[v,0]\to T(v)$. Now deriving
\[T(1),\ldots,T(v-1)\to V[v,0],T(v)\tag{4.1}\]
will allow us to derive $T(1),\ldots,T(v-1)\to T(v)$ by cutting on $V[v,0]$. We can derive(4.1) from the formulas
\[T(1),\ldots,T(v-1),V[v,i]\to T(v)\tag{4.2}\]
for $i\in[n]$ using the clauses(3.11) for $V[v,i]$. The formulas(4.2) can be in turn derived from the formulas
\[T(v_\ell),L[v,v_\ell],T(v_r),R[u,v_r],V[v,i]\to T(v)\tag{4.3}\]
for $v_\ell,v_r\in[s]$ using the clauses(3.11) for $L[v,k]$ and $R[v,k]$,(3.12) and(3.13). Finally,(4.3) can be derived from the formulas
\[D[v_\ell,j,\alpha(x_j),L[v,v_\ell],D[v_r,x_k,\alpha(x_k)],R[u,v_r],V[v,i]\to T(v),\tag{4.4}\]
for $j,k\in[n]$, which can be derived directly from the clauses(3.21) and either(3.14),(3.15) and(3.19) or(3.16) and(3.17) depending on whether $i=j=k$ or not.
We can see that the derivations of $T(1),\overline{T(s)}$ and $T(1),\ldots,T(v-1)\to T(v)$ take at most $O(m+s^2n^3)$ steps, hence the overall refutation has size $O(s^2n(m+s^2n^3))$.
Now, notice that after substituting every variable $P[u]$ in it with $S_{d,N}^{\wedge}(P[u])$ and every other variable $x$ with $S_{d,N}^{\vee}(x)$, $T(v)$ becomes a depth-$d$ formula. Hence we see that after the substitution, the refutation described above becomes a depth-$(d+1)$ tree-like LK refutation of $\RRef_{d,N}(F,s)$. We can then get a depth-$d$ refutation of $\RRef_{d,N}(F,s)$ of the required size by applying Proposition2.2.
\end{proof}
